\editionPage[\pdfbookmark[0]{Défis d’approfondissement}{bm-defis-debut-parent}]{defis-debut}{Dix-huit défis}{Un défi par chapitre · Parcours facultatif}{Chercher avant de consulter les deux indications.}
\begin{blocEdition}{defis-debut}
Ces problèmes demandent davantage d’initiative. Ils utilisent les outils du livre ; les prolongements sont définis dans l’énoncé. Les défis sont facultatifs et approfondissent les outils du programme.\par\bigskip
\begin{savoirfaire}[Une recherche qui avance]Faire un schéma, tester un petit cas, comparer deux méthodes, conserver les essais et préciser ce qui reste à prouver. Ouvrir la première indication seulement après une recherche réelle ; revenir au problème avant de lire la seconde.\end{savoirfaire}
\begin{savoirfaire}[Après la solution]Identifier l’idée décisive et la refaire sur une feuille vierge. Expliquer pourquoi elle fonctionne compte autant que retrouver le résultat.\end{savoirfaire}
\end{blocEdition}
\editionPage{defi-ex-1}{D01 · Défi}{Cent termes, sans cent calculs}{Approfondissement · 30 à 60 minutes indicatives}
\begin{blocEdition}{defi-ex-1}
Calculer $S=-1+2-3+4-\cdots-99+100$ sans additionner terme par terme. Expliquer le résultat. Que vaut la somme si l’on s’arrête à $-99$ ?\par\vspace{5mm}{\otLight\fontsize{8}{10}\selectfont\color{otMuted}Espace de recherche}\par\par\vspace{7mm}{\color{otRule!65}\hrule height .25pt}\par\vspace{7mm}{\color{otRule!65}\hrule height .25pt}\par\vspace{7mm}{\color{otRule!65}\hrule height .25pt}\par\vspace{7mm}{\color{otRule!65}\hrule height .25pt}\par\vspace{7mm}{\color{otRule!65}\hrule height .25pt}\par\vspace{7mm}{\color{otRule!65}\hrule height .25pt}
\end{blocEdition}
\editionPage{defi-ex-2}{D02 · Défi}{Retrouver tout le budget}{Approfondissement · 30 à 60 minutes indicatives}
\begin{blocEdition}{defi-ex-2}
On utilise un tiers d’un budget, puis un quart du reste, puis trois cinquièmes du nouveau reste. Il reste alors 18 euros. Retrouver le budget et les trois sommes utilisées.\par\vspace{5mm}{\otLight\fontsize{8}{10}\selectfont\color{otMuted}Espace de recherche}\par\par\vspace{7mm}{\color{otRule!65}\hrule height .25pt}\par\vspace{7mm}{\color{otRule!65}\hrule height .25pt}\par\vspace{7mm}{\color{otRule!65}\hrule height .25pt}\par\vspace{7mm}{\color{otRule!65}\hrule height .25pt}\par\vspace{7mm}{\color{otRule!65}\hrule height .25pt}\par\vspace{7mm}{\color{otRule!65}\hrule height .25pt}
\end{blocEdition}
\editionPage{defi-ex-3}{D03 · Défi}{Jusqu’où peut-on doubler ?}{Approfondissement · 30 à 60 minutes indicatives}
\begin{blocEdition}{defi-ex-3}
Une mémoire fictive contient au départ un seul élément. À chaque étape, le nombre d’éléments double. Sa capacité maximale est un million d’éléments. Quel est le plus grand nombre d’étapes possible ? Donner un encadrement avec deux puissances de 2.\par\vspace{5mm}{\otLight\fontsize{8}{10}\selectfont\color{otMuted}Espace de recherche}\par\par\vspace{7mm}{\color{otRule!65}\hrule height .25pt}\par\vspace{7mm}{\color{otRule!65}\hrule height .25pt}\par\vspace{7mm}{\color{otRule!65}\hrule height .25pt}\par\vspace{7mm}{\color{otRule!65}\hrule height .25pt}\par\vspace{7mm}{\color{otRule!65}\hrule height .25pt}\par\vspace{7mm}{\color{otRule!65}\hrule height .25pt}
\end{blocEdition}
\editionPage{defi-ex-4}{D04 · Défi}{Quatre carreaux pour un grand carré}{Approfondissement · 30 à 60 minutes indicatives}
\begin{blocEdition}{defi-ex-4}
Quatre carreaux carrés identiques, chacun d’aire 18 cm$^2$, sont assemblés en un grand carré, en deux lignes et deux colonnes. Déterminer le côté et le périmètre exacts du grand carré. Comparer le rapport des côtés et celui des aires.\label{fig-defi-4}\figureGeo{defi-racine}
\par\vspace{5mm}{\otLight\fontsize{8}{10}\selectfont\color{otMuted}Espace de recherche}\par\par\vspace{7mm}{\color{otRule!65}\hrule height .25pt}\par\vspace{7mm}{\color{otRule!65}\hrule height .25pt}\par\vspace{7mm}{\color{otRule!65}\hrule height .25pt}\par\vspace{7mm}{\color{otRule!65}\hrule height .25pt}\par\vspace{7mm}{\color{otRule!65}\hrule height .25pt}\par\vspace{7mm}{\color{otRule!65}\hrule height .25pt}
\end{blocEdition}
\editionPage{defi-ex-5}{D05 · Défi}{Comparer sans choisir de valeur}{Approfondissement · 30 à 60 minutes indicatives}
\begin{blocEdition}{defi-ex-5}
Pour $x>0$, un rectangle A mesure 3 cm sur $x+4$ cm. Un rectangle B mesure 2 cm sur $x+6$ cm. Quelle aire est la plus grande ? Est-il possible qu’elles soient égales pour une longueur $x$ strictement positive ?\label{fig-defi-5}\figureGeo{defi-litteral}
\par\vspace{5mm}{\otLight\fontsize{8}{10}\selectfont\color{otMuted}Espace de recherche}\par\par\vspace{7mm}{\color{otRule!65}\hrule height .25pt}\par\vspace{7mm}{\color{otRule!65}\hrule height .25pt}\par\vspace{7mm}{\color{otRule!65}\hrule height .25pt}\par\vspace{7mm}{\color{otRule!65}\hrule height .25pt}\par\vspace{7mm}{\color{otRule!65}\hrule height .25pt}\par\vspace{7mm}{\color{otRule!65}\hrule height .25pt}
\end{blocEdition}
\editionPage{defi-ex-6}{D06 · Défi}{Compter les pièces}{Approfondissement · 30 à 60 minutes indicatives}
\begin{blocEdition}{defi-ex-6}
Une boîte contient 42 pièces, uniquement de 1 euro et de 2 euros. La somme vaut 68 euros. Trouver le nombre de pièces de chaque sorte. Justifier que les deux contraintes sont satisfaites.\par\vspace{5mm}{\otLight\fontsize{8}{10}\selectfont\color{otMuted}Espace de recherche}\par\par\vspace{7mm}{\color{otRule!65}\hrule height .25pt}\par\vspace{7mm}{\color{otRule!65}\hrule height .25pt}\par\vspace{7mm}{\color{otRule!65}\hrule height .25pt}\par\vspace{7mm}{\color{otRule!65}\hrule height .25pt}\par\vspace{7mm}{\color{otRule!65}\hrule height .25pt}\par\vspace{7mm}{\color{otRule!65}\hrule height .25pt}
\end{blocEdition}
\editionPage{defi-ex-7}{D07 · Défi}{Annuler une remise}{Approfondissement · 30 à 60 minutes indicatives}
\begin{blocEdition}{defi-ex-7}
Un prix baisse de 20\,\%. De quel pourcentage faut-il augmenter le nouveau prix pour retrouver l’ancien ? Prouver que le pourcentage obtenu ne dépend pas du prix initial positif.\par\vspace{5mm}{\otLight\fontsize{8}{10}\selectfont\color{otMuted}Espace de recherche}\par\par\vspace{7mm}{\color{otRule!65}\hrule height .25pt}\par\vspace{7mm}{\color{otRule!65}\hrule height .25pt}\par\vspace{7mm}{\color{otRule!65}\hrule height .25pt}\par\vspace{7mm}{\color{otRule!65}\hrule height .25pt}\par\vspace{7mm}{\color{otRule!65}\hrule height .25pt}\par\vspace{7mm}{\color{otRule!65}\hrule height .25pt}
\end{blocEdition}
\editionPage{defi-ex-8}{D08 · Défi}{Un remplissage avec fuite}{Approfondissement · 30 à 60 minutes indicatives}
\begin{blocEdition}{defi-ex-8}
Un bassin de 900 L contient 75 L au départ. Un robinet débite 18 L/min ; une fuite retire 3 L/min. Après 20 min, la fuite est réparée et le robinet continue. Combien de temps faut-il au total pour remplir le bassin ?\par\vspace{5mm}{\otLight\fontsize{8}{10}\selectfont\color{otMuted}Espace de recherche}\par\par\vspace{7mm}{\color{otRule!65}\hrule height .25pt}\par\vspace{7mm}{\color{otRule!65}\hrule height .25pt}\par\vspace{7mm}{\color{otRule!65}\hrule height .25pt}\par\vspace{7mm}{\color{otRule!65}\hrule height .25pt}\par\vspace{7mm}{\color{otRule!65}\hrule height .25pt}\par\vspace{7mm}{\color{otRule!65}\hrule height .25pt}
\end{blocEdition}
\editionPage{defi-ex-9}{D09 · Défi}{Un découpage variable}{Approfondissement · 30 à 60 minutes indicatives}
\begin{blocEdition}{defi-ex-9}
Dans un rectangle de 18 cm sur 10 cm, on retire quatre carrés de côté $x$ cm aux coins, avec $0<x<5$. Donner la formule de l’aire restante. Parmi $x=2,3,4$, lequel laisse 116 cm$^2$ ? Peut-on retrouver ce choix sans tester les trois ?\label{fig-defi-9}\figureGeo{coins-retires}
\par\vspace{5mm}{\otLight\fontsize{8}{10}\selectfont\color{otMuted}Espace de recherche}\par\par\vspace{7mm}{\color{otRule!65}\hrule height .25pt}\par\vspace{7mm}{\color{otRule!65}\hrule height .25pt}\par\vspace{7mm}{\color{otRule!65}\hrule height .25pt}\par\vspace{7mm}{\color{otRule!65}\hrule height .25pt}\par\vspace{7mm}{\color{otRule!65}\hrule height .25pt}\par\vspace{7mm}{\color{otRule!65}\hrule height .25pt}
\end{blocEdition}
\editionPage{defi-ex-10}{D10 · Défi}{Même moyenne, autre médiane}{Approfondissement · 30 à 60 minutes indicatives}
\begin{blocEdition}{defi-ex-10}
Comparer les séries A : 4, 4, 4, 16, 16, 16 et B : 0, 8, 8, 8, 8, 28. Calculer moyenne et médiane. Que peut-on conclure sur l’information fournie par une moyenne seule ?\par\vspace{5mm}{\otLight\fontsize{8}{10}\selectfont\color{otMuted}Espace de recherche}\par\par\vspace{7mm}{\color{otRule!65}\hrule height .25pt}\par\vspace{7mm}{\color{otRule!65}\hrule height .25pt}\par\vspace{7mm}{\color{otRule!65}\hrule height .25pt}\par\vspace{7mm}{\color{otRule!65}\hrule height .25pt}\par\vspace{7mm}{\color{otRule!65}\hrule height .25pt}\par\vspace{7mm}{\color{otRule!65}\hrule height .25pt}
\end{blocEdition}
\editionPage{defi-ex-11}{D11 · Défi}{Une cible et une comparaison exacte}{Approfondissement · 30 à 60 minutes indicatives}
\begin{blocEdition}{defi-ex-11}
Un point est choisi uniformément dans un carré de côté 8 cm. On gagne s’il appartient au disque central de rayon 2 cm. Donner la probabilité exacte et montrer qu’elle est inférieure à $1/5$, en utilisant $\pi<3{,}2$.\label{fig-defi-11}\figureGeo{defi-cible}
\par\vspace{5mm}{\otLight\fontsize{8}{10}\selectfont\color{otMuted}Espace de recherche}\par\par\vspace{7mm}{\color{otRule!65}\hrule height .25pt}\par\vspace{7mm}{\color{otRule!65}\hrule height .25pt}\par\vspace{7mm}{\color{otRule!65}\hrule height .25pt}\par\vspace{7mm}{\color{otRule!65}\hrule height .25pt}\par\vspace{7mm}{\color{otRule!65}\hrule height .25pt}\par\vspace{7mm}{\color{otRule!65}\hrule height .25pt}
\end{blocEdition}
\editionPage{defi-ex-12}{D12 · Défi}{Une droite, trois rôles}{Approfondissement · 30 à 60 minutes indicatives}
\begin{blocEdition}{defi-ex-12}
ABC est isocèle en A et M est le milieu de BC. En comparant deux triangles, montrer que AM est à la fois une médiane, une hauteur et la bissectrice de l’angle A. Une bissectrice partage un angle en deux angles égaux.\label{fig-defi-12}\figureGeo{defi-triangles}
\par\vspace{5mm}{\otLight\fontsize{8}{10}\selectfont\color{otMuted}Espace de recherche}\par\par\vspace{7mm}{\color{otRule!65}\hrule height .25pt}\par\vspace{7mm}{\color{otRule!65}\hrule height .25pt}\par\vspace{7mm}{\color{otRule!65}\hrule height .25pt}\par\vspace{7mm}{\color{otRule!65}\hrule height .25pt}\par\vspace{7mm}{\color{otRule!65}\hrule height .25pt}\par\vspace{7mm}{\color{otRule!65}\hrule height .25pt}
\end{blocEdition}
\editionPage{defi-ex-13}{D13 · Défi}{Le câble et ses attaches}{Approfondissement · 30 à 60 minutes indicatives}
\begin{blocEdition}{defi-ex-13}
Un câble doit suivre l’hypoténuse d’un triangle rectangle de côtés perpendiculaires 9 m et 12 m. Il faut en plus 0,75 m à chaque extrémité pour les attaches. Le câble coûte 4,80 euros par mètre. Calculer la longueur et le coût exacts.\label{fig-defi-13}\figureGeo{defi-pyth}
\par\vspace{5mm}{\otLight\fontsize{8}{10}\selectfont\color{otMuted}Espace de recherche}\par\par\vspace{7mm}{\color{otRule!65}\hrule height .25pt}\par\vspace{7mm}{\color{otRule!65}\hrule height .25pt}\par\vspace{7mm}{\color{otRule!65}\hrule height .25pt}\par\vspace{7mm}{\color{otRule!65}\hrule height .25pt}\par\vspace{7mm}{\color{otRule!65}\hrule height .25pt}\par\vspace{7mm}{\color{otRule!65}\hrule height .25pt}
\end{blocEdition}
\editionPage{defi-ex-14}{D14 · Défi}{Une bande entre triangles}{Approfondissement · 30 à 60 minutes indicatives}
\begin{blocEdition}{defi-ex-14}
ABC est rectangle en A, $AB=12$ et $AC=16$ cm. M est sur AB, avec $AM=3$ cm. La parallèle à BC par M coupe AC en N. Calculer AN, MN, puis l’aire de la partie de ABC extérieure à AMN.\label{fig-defi-14}\figureGeo{defi-thales}
\par\vspace{5mm}{\otLight\fontsize{8}{10}\selectfont\color{otMuted}Espace de recherche}\par\par\vspace{7mm}{\color{otRule!65}\hrule height .25pt}\par\vspace{7mm}{\color{otRule!65}\hrule height .25pt}\par\vspace{7mm}{\color{otRule!65}\hrule height .25pt}\par\vspace{7mm}{\color{otRule!65}\hrule height .25pt}\par\vspace{7mm}{\color{otRule!65}\hrule height .25pt}\par\vspace{7mm}{\color{otRule!65}\hrule height .25pt}
\end{blocEdition}
\editionPage{defi-ex-15}{D15 · Défi}{Choisir la bonne inclinaison}{Approfondissement · 30 à 60 minutes indicatives}
\begin{blocEdition}{defi-ex-15}
Une barre de 10 m forme l’hypoténuse d’un triangle rectangle. Son côté horizontal adjacent à l’angle étudié peut mesurer 5, 6 ou 7 m. Quel choix donne un angle compris entre $50^\circ$ et $55^\circ$ ? Justifier les trois essais.\label{fig-defi-15}\figureGeo{defi-cosinus}
\par\vspace{5mm}{\otLight\fontsize{8}{10}\selectfont\color{otMuted}Espace de recherche}\par\par\vspace{7mm}{\color{otRule!65}\hrule height .25pt}\par\vspace{7mm}{\color{otRule!65}\hrule height .25pt}\par\vspace{7mm}{\color{otRule!65}\hrule height .25pt}\par\vspace{7mm}{\color{otRule!65}\hrule height .25pt}\par\vspace{7mm}{\color{otRule!65}\hrule height .25pt}\par\vspace{7mm}{\color{otRule!65}\hrule height .25pt}
\end{blocEdition}
\editionPage{defi-ex-16}{D16 · Défi}{Deux glissements, une construction}{Approfondissement · 30 à 60 minutes indicatives}
\begin{blocEdition}{defi-ex-16}
Le triangle a pour sommets A(0;1), B(2;1), C(1;3). On le déplace de 4 carreaux à droite et 2 en bas, puis de 1 à gauche et 5 en haut. Donner les deux triangles images. Le déplacement final conserve-t-il l’aire ?\label{fig-defi-16}\figureGeo{defi-translation}
\par\vspace{5mm}{\otLight\fontsize{8}{10}\selectfont\color{otMuted}Espace de recherche}\par\par\vspace{7mm}{\color{otRule!65}\hrule height .25pt}\par\vspace{7mm}{\color{otRule!65}\hrule height .25pt}\par\vspace{7mm}{\color{otRule!65}\hrule height .25pt}\par\vspace{7mm}{\color{otRule!65}\hrule height .25pt}\par\vspace{7mm}{\color{otRule!65}\hrule height .25pt}\par\vspace{7mm}{\color{otRule!65}\hrule height .25pt}
\end{blocEdition}
\editionPage{defi-ex-17}{D17 · Défi}{Un volume sous deux formes}{Approfondissement · 30 à 60 minutes indicatives}
\begin{blocEdition}{defi-ex-17}
Une pyramide de base carrée de côté 6 cm et de hauteur 8 cm est comparée à un prisme droit de même base et de même hauteur. Calculer les deux volumes. Combien de remplissages de la pyramide faut-il verser pour remplir le prisme sans perte ? Puis réduire les deux solides de moitié.\label{fig-defi-17}\figureGeo{defi-pyramide}
\par\vspace{5mm}{\otLight\fontsize{8}{10}\selectfont\color{otMuted}Espace de recherche}\par\par\vspace{7mm}{\color{otRule!65}\hrule height .25pt}\par\vspace{7mm}{\color{otRule!65}\hrule height .25pt}\par\vspace{7mm}{\color{otRule!65}\hrule height .25pt}\par\vspace{7mm}{\color{otRule!65}\hrule height .25pt}\par\vspace{7mm}{\color{otRule!65}\hrule height .25pt}\par\vspace{7mm}{\color{otRule!65}\hrule height .25pt}
\end{blocEdition}
\editionPage{defi-ex-18}{D18 · Défi}{Une somme programmée et prouvée}{Approfondissement · 30 à 60 minutes indicatives}
\begin{blocEdition}{defi-ex-18}
Exécuter mentalement le script pour les entiers de 1 à 10. Donner s et n à la fin. Retrouver la somme en regroupant les entiers par paires. Que se passerait-il si « mettre s à 0 » était déplacé à l’intérieur de la boucle ?\label{fig-defi-18}\figureGeo{defi-script}
\par\vspace{5mm}{\otLight\fontsize{8}{10}\selectfont\color{otMuted}Espace de recherche}\par\par\vspace{7mm}{\color{otRule!65}\hrule height .25pt}\par\vspace{7mm}{\color{otRule!65}\hrule height .25pt}\par\vspace{7mm}{\color{otRule!65}\hrule height .25pt}\par\vspace{7mm}{\color{otRule!65}\hrule height .25pt}\par\vspace{7mm}{\color{otRule!65}\hrule height .25pt}\par\vspace{7mm}{\color{otRule!65}\hrule height .25pt}
\end{blocEdition}
